\documentclass[11pt]{beamer}
\usetheme{Madrid}
\usepackage[spanish]{babel}
\usepackage[utf8]{inputenc}
\usepackage{amsmath}
\usepackage{amsfonts}
\usepackage{amssymb}
\usepackage{graphicx}
\author{Elena María Gómez Ríos y Jose Luis Martínez Ortiz}
\title{Centralitas de voz virtual} % (asterisk, 3CX, telefacil.com)
%\setbeamercovered{transparent} 
%\setbeamertemplate{navigation symbols}{} 
%\logo{} 
%\institute{} 
%\date{} 
\subject{Fundamentos de Redes} 
\begin{document}

\begin{frame}
\titlepage
\end{frame}

%\begin{frame}
%\tableofcontents
%\end{frame}

\begin{frame}{Apartados}
\begin{enumerate}
\item Introducción motivadora
\item Aplicaciones y plataformas existentes.
\item Ejemplo de uso.
\item Modelo de negocio y/o explotación.
\item Protocolos.
\end{enumerate}
\end{frame}


\begin{frame}{Introducción motivadora}
% http://www.voz.com/centralita-virtual.php
\begin{enumerate}[]
\item \begin{center}
¿Qué son las Centralitas de Voz?
\end{center}
\item 
\item \begin{center}
Son un \textbf{sistema telefónico avanzado} por el cual se puede obtener tantas extensiones como se desee, fijas o móviles, de forma virtual. 
\\Una extensión es cada puesto de teléfono.
\end{center}


\end{enumerate}
\end{frame}


\begin{frame}{Introducción motivadora}
\textbf{¿Qué ventajas tiene frente a una centralita física?}\\
\vspace{1cm}

\begin{enumerate}[$\bullet$]
\item \textbf{No} se queda \textbf{obsoleta} con el tiempo.
\item Utiliza la red de \textbf{Internet}, pero depende del ancho de banda contratado.
\item Utiliza canales de voz virtuales, por lo tanto no están limitados. 
% Es más se pueden crear tantos canales de voz internos como se necesite.
\item \textbf{No} tiene \textbf{coste} de establecimiento de llamada en la red interna.
\item Es un \textbf{ahorro} considerable con respecto a la telefonía tradicional.
\end{enumerate}
\end{frame} 

\begin{frame}{Distribuidores de Centralitas}
Algunas compañías de servicio de centralitas nacionales son:
\begin{enumerate}
\item netelip
\item VOZ.com
\item Megacall
\end{enumerate}
\vspace*{1cm}
En el ámbito internacional están:
\begin{enumerate}
\item Digium (The Asterisk Company)
\item 3CX
\item AVOXI
\end{enumerate}

\end{frame}


\begin{frame}{Ejemplos de uso}
\begin{figure}[H] %con el [H] le obligamos a situar aquí la figura
\centering
\includegraphics[scale=0.4]{./imagenes/logos.png}
\end{figure}

\end{frame}


\begin{frame}{Modelos de Negocio}
Principalmente el modelo de negocio es la \textbf{venta de centralitas de voz 
ya montadas} y funcionando a empresas y entidades públicas por un coste 
mensual dependiendo de la amplitud de la centralita, además del hardware
necesario para su utilización como teléfonos, switches, etc.
\\
\vspace*{1cm}
Aunque ``Asterisk'' se centra en dar un software gratuito para que terceros vendan centralitas a partir de éste. Su modelo de negocio es el de donaciones y patrocinadores.
\end{frame}



\begin{frame}{Protocolos VOIP}
Para llevar a cabo la transmisión en tiempo real de una comunicación 
utiliza un protocolo de tipo \textbf{VOIP} (Voice over IP), algunos más
famosos son:
\vspace*{0.5cm}
\begin{enumerate}
\item Skype, propiedad de \texttt{Microsoft}
\item IAX2, original de \texttt{Asterisk} y libre
\item SCCP, propiedad de \texttt{CISCO}
\end{enumerate}
\end{frame}

\begin{frame}{Protocolo IAX2}
Utiliza UDP y está orientado al streaming media, pero fue diseñado principalmente para llamadas de voz por IP.\\
\vspace*{1cm}

Por defecto utiliza el puerto 4569 y transmite los datos por ``in-band'' 
lo que hace que el protocolo sea muy transparente para el cortafuegos y
 así ser muy rápido en redes internas.\\
\vspace*{0.5cm}
Soporta ``\texttt{Trunking red}'', lo que permite enviar varias llamadas a la vez en un único paquete de UDP (datagram), con el ahorro de latencia que esto supone. \\
\vspace*{0.5cm}
+información en el RFC 5456
\end{frame}

\begin{frame}{Protocolo Skype}
Skype utiliza el protocolo de transporte UDP para la transmisión de datos y TCP para la identificación de los usuarios. \\
\vspace*{0.3cm}
Por defecto utiliza el puerto 1234.\\
\vspace*{0.3cm}
Skype utiliza el modelo P2P en vez del modelo cliente-servidor.

\vspace*{1cm}
+información en el RFC 5688
\end{frame}

\begin{frame}{Protocolo MSNP24}
Microsoft Notification Protocol. \\
\vspace*{0.3cm}
Este protocolo ha sido desarrollado por Microsoft para sus servicios de mensajería, como Skype (sustituyendo el protocolo Skype).\\
\vspace*{0.3cm}
Sobre este protocolo no hay demasiada información al ser privativo y actual.
\end{frame}


\end{document}